\documentclass[12pt, a4paper]{article}

% --- Pacotes Básicos ---
\usepackage[utf8]{inputenc}
\usepackage[T1]{fontenc}
%\usepackage[main=brazilian, english]{babel}
\usepackage{indentfirst}
\usepackage{geometry}
\geometry{left=3cm, top=3cm, right=2cm, bottom=2cm}

% --- Informações do Documento ---
\title{Resenha}
\author{Erickson Giesel Müller}
\date{12 de Maio de 2026}

\begin{document}
\maketitle
\thispagestyle{empty} % primeira página não tem número.

O trabalho intitulado "CONTINUIDADE E MUDANÇA DO TERRITÓRIO KAINGANG DO RIO GRANDE DO SUL: um estudo de caso do aldeamento de Nonoai" constitui-se como o Trabalho de Conclusão de Curso em História da autora Jussara Costa Curta, apresentado na Universidade Federal do Rio Grande do Sul (UFRGS) no ano de 2012. Sob a orientação da Profa. Dra. Adriana Schmidt Dias, doutora em arqueologia pela Universidade de São Paulo (USP) e professora titular do Departamento de História da UFRGS. O estudo estabelece como objetivo central a análise do processo de expropriação e perda das terras tradicionais Kaingang na Província do Rio Grande do Sul, focando-se especificamente no marco temporal de 1848 a 1889. Além da reconstrução histórica através do estudo de caso de Nonoai, a pesquisa busca compreender as concepções nativas de territorialidade e identificar as múltiplas estratégias de resistência e negociação utilizadas por esse povo indígena para a manutenção de seu espaço frente ao avanço da colonização.  

A delimitação espacial da pesquisa concentra-se no norte do estado do Rio Grande do Sul, tendo como foco central o estudo de caso do aldeamento de Nonoai. O recorte temporal abrange o período de 1848 a 1889, marco que se inicia com o processo de formação do aldeamento e se estende até o encerramento do período Imperial no Brasil. O objetivo fundamental dessa delimitação é investigar o processo de expropriação do território nativo Kaingang e analisar as diferentes estratégias, tanto de resistência quanto de negociação, mobilizadas pelos indígenas para a manutenção de suas terras tradicionais diante do avanço da colonização.

O propósito central do estudo é analisar o processo de perda e expropriação do território tradicional Kaingang na Província do Rio Grande do Sul. Paralelamente à reconstrução histórica desse processo, a pesquisa visa conhecer a concepção nativa de território e identificar as estratégias utilizadas pelos indígenas para a manutenção e defesa de suas terras. O foco no aldeamento de Nonoai justifica-se pelo fato de este núcleo ter sido um dos que mais promoveu ações e reações frente à expropriação territorial, permitindo compreender a agência indígena e sua resistência diante do contexto colonizador.

A filiação linguística dos Kaingang ao Tronco Macro-Jê e à Família Jê revela uma trajetória histórica que remonta originalmente ao Brasil Central, em áreas situadas entre as bacias dos rios São Francisco, Araguaia e Tocantins. De acordo com modelos de reconstrução linguística, esse grupo teria se separado de outros povos Jê há cerca de 3 mil anos, iniciando um longo processo migratório rumo ao sul do continente. Essa movimentação não foi aleatória: os grupos buscaram regiões de planalto no Sul do Brasil que fossem geomorfologicamente semelhantes ao seu habitat de origem, permitindo a manutenção de padrões de adaptação e subsistência já consolidados.  

A arqueologia fundamenta a tese da continuidade histórica dos Kaingang através da chamada "Tradição Taquara", identificada pela presença de pequenos vasos cerâmicos e estruturas de terra nos sítios arqueológicos. O elemento mais emblemático dessa tradição são as "casas subterrâneas", moradias escavadas em regiões de altitude acima de 500 metros que serviam como uma proteção engenhosa contra os invernos rigorosos do Planalto Meridional. As pesquisas demonstram que as áreas dessas ocupações pré-históricas, que remontam ao século II d.C., coincidem geograficamente com os territórios ocupados pelos Kaingang em períodos posteriores, comprovando raízes culturais milenares na região. A evidência definitiva dessa permanência é reforçada por datações de Carbono 14, que mostram a utilização dessas estruturas até meados do século XIX, período em que registros históricos de viajantes e engenheiros já descreviam a presença desses mesmos grupos habitando as referidas áreas.

A ocupação histórica dos Kaingang no Sul do Brasil demonstra uma relação de profunda integração com o meio ambiente, coincidindo quase perfeitamente com a distribuição das florestas de araucária. Essas matas eram vitais para a subsistência do grupo, fornecendo o pinhão, um recurso alimentar estratégico que servia tanto para a nutrição imediata quanto para provisões durante as estações de escassez. Além de sustentar a dieta, o ambiente florestal das terras altas oferecia proteção essencial contra o clima frio e os ventos das grandes altitudes, influenciando inclusive o modelo de construção das moradias do grupo. O vínculo com esse ecossistema era tão expressivo que os pinheirais eram rigorosamente divididos entre as tribos, com limites demarcados nas cascas dos pinheiros através de marcas e pinturas que refletiam a identidade social e o direito de uso territorial dos ocupantes.

No século XIX, as denominações atribuídas aos grupos Jê meridionais eram variadas e frequentemente baseadas em observações externas dos colonizadores. Seus ancestrais diretos eram historicamente identificados como Guaianá, termo que englobava diversas tribos com costumes e línguas distintos dos Guarani. Entretanto, até meados do século XIX, a denominação que se tornou generalizada na Província do Rio Grande do Sul e em outras partes do Império foi a de "Coroados". De acordo com os registros do engenheiro Alphonse Mabilde(1983), esse nome derivava do costume desses indígenas de tonsurarem seus cabelos. O termo "Kaingang" só começou a surgir na documentação bibliográfica a partir de 1882, sendo consolidado posteriormente como uma designação genérica para os diversos grupos de dialetos aparentados e filiados ao tronco linguístico Jê.

A estrutura política dos Kaingang no século XIX era organizada em pequenas tribos compostas por núcleos familiares e parentes próximos. Cada uma dessas unidades possuía seu próprio cacique local, o qual respondia hierarquicamente a um cacique principal. Este líder central exercia autoridade direta sobre a gestão do território, sendo responsável por indicar o local específico das matas de pinheiros que cada grupo subordinado deveria habitar. A estabilidade dessas relações dependia de um complexo sistema de visitas periódicas de caráter político; quando essas visitas eram suspensas por motivo de insubordinação ou desconfiança, o ato era interpretado como uma declaração formal de guerra. Esse cenário era frequentemente marcado por uma dinâmica de conflitos constantes, envolvendo tanto dissidências internas quanto disputas territoriais históricas contra grupos rivais, como os Xokleng e os Guaranis.

A economia de subsistência dos Kaingang no século XIX era fundamentada em um sistema diversificado e profundamente integrado aos ciclos da natureza, garantindo a sobrevivência do grupo por meio da exploração de diferentes ecossistemas. A agricultura era praticada em terrenos elevados, onde se cultivavam diversas variedades de milho, feijão, abóbora e amendoim. A coleta desempenhava um papel central na dieta e na cultura, com destaque absoluto para o pinhão, o alimento mais prestigiado e cuja exploração seguia regras rígidas de propriedade territorial. Além do pinhão, o grupo coletava mel, palmito, larvas de taquara (corós), raízes e uma grande variedade de frutos sazonais, como o araçá, a pitanga e a guabiroba. A obtenção de proteína animal era assegurada pela caça de espécies como a anta, o porco do mato, o tatu e a paca em florestas fechadas, além da pesca realizada durante os meses de verão em locais mais afastados das aldeias. Todo esse sistema era regido por uma dinâmica de mobilidade sazonal: o ano iniciava com o cultivo das roças na primavera, seguido por expedições de pesca no verão e a colheita do pinhão no outono, permitindo que as comunidades vivessem de provisões estocadas durante o rigoroso inverno do planalto.

Para os Kaingang, a terra transcende a função de mero recurso de subsistência, constituindo o suporte fundamental da vida social e o pilar central de seu sistema de crenças e conhecimento. Essa concepção de territorialidade possui uma dimensão sociopolítica e cosmológica vasta, na qual a ligação com a "terra-mãe" é estabelecida desde o nascimento e mantida até a morte. O território é, portanto, um espaço carregado de conteúdo simbólico e histórico, estruturado por explicações míticas fundamentais, como a do grande dilúvio, no qual os antepassados Kaingang sobreviveram ao buscar o cume da serra Krinjijimbé (Serra do Mar), que sobressaía das águas. Diferente da visão puramente geográfica e utilitária da sociedade nacional, para o indígena a relação com o território tradicional é vital para a afirmação de sua identidade étnica, integrando as dimensões materiais de caça e coleta à preservação da memória e ao apego aos locais onde seus mortos estão sepultados.

A noção de propriedade para a sociedade Kaingang difere fundamentalmente do conceito capitalista de propriedade privada, sendo pautada por relações sociais e pelo direito de uso coletivo do território. Ter uma propriedade comum do solo não significava a ausência de normas; pelo contrário, existiam direitos precisos e definidos que regulavam o acesso ao espaço. Enquanto as florestas eram espaços de caça e coleta acessíveis a qualquer indivíduo do grupo, os pinheirais possuíam uma lógica específica de divisão, na qual cada subgrupo detinha uma parcela determinada para a colheita do pinhão. Esse sistema de direitos comunitários, longe de ser uma restrição, assegurava que todos os membros tivessem acesso aos recursos fundamentais e, crucialmente, garantia que as gerações futuras recebessem a herança desse território preservado.

O texto promove uma desconstrução crítica de mitos historiográficos utilizados pelo Estado para legitimar o deslocamento dos povos indígenas, especificamente as noções de "nomadismo" e de "vazio demográfico". A ideia de nomadismo é caracterizada como uma construção ideológica destinada a justificar a expropriação de terras e a destituição dos direitos nativos sobre o solo. Contrapondo-se a essa visão, autores citados no estudo argumentam que os Kaingang já eram agricultores antes do contato com os europeus e que sua estratégia de circulação sazonal pelo território para coleta e caça foi erroneamente interpretada como uma ausência de morada fixa. Paralelamente, o conceito de "vazio demográfico" é apresentado como um discurso oficial que visava ocultar a presença indígena, tratando as terras do planalto como "devolutas" ou desabitadas para viabilizar uma ocupação considerada "produtiva" pela lógica colonial. Essa narrativa serviu para apresentar a colonização europeia como um processo pacífico em terras sem dono, ignorando as evidências arqueológicas de ocupação milenar e os registros históricos de resistência ativa dos grupos indígenas contra a invasão de seus domínios tradicionais.

A Lei de Terras de 1850 institucionalizou a exclusão das comunidades indígenas ao estabelecer novos princípios para a ocupação do solo que favoreciam a mercantilização da terra. A legislação exigia, em seus artigos 5º e 6º, a comprovação de "morada habitual" e "cultura efetiva" para a legitimação das posses, parâmetros que ignoravam o manejo territorial e a lógica de uso do solo das sociedades nativas. Como as atividades fundamentais dos Kaingang como a caça, a pesca, a coleta sazonal e a mobilidade geográfica não se enquadravam nas exigências de vida sedentária e agricultura permanente da nova lei, seus roçados e habitações tradicionais não foram considerados ocupações legítimas. Dessa forma, vastas áreas de domínio indígena foram sistematicamente classificadas como "terras devolutas", sendo disponibilizadas pelo Estado para a expansão do latifúndio e para os projetos de colonização europeia, privando os indígenas do acesso legal aos seus territórios ancestrais.  

% A Lei de Terras inaugurou uma política agressiva de mercantilização do solo ao instituir que o acesso às terras públicas passaria a ocorrer exclusivamente por meio da compra, transformando o território em uma mercadoria central para a expansão da economia agrícola capitalista. Essa legislação criou o mecanismo jurídico necessário para reclassificar vastas áreas tradicionalmente ocupadas como "terras devolutas" — terrenos tidos como desocupados ou sem título legítimo — facilitando sua apropriação por grandes latifundiários que desejavam ampliar suas posses. Ao proteger os interesses das elites fundiárias e disponibilizar essas áreas para o projeto de colonização europeia, o Estado garantiu a reserva de terras para o imigrante considerado "produtivo", enquanto promovia a retirada sistemática dos indígenas de seus domínios originais sob a justificativa de que estes não constituíam ocupantes efetivos.  

O aldeamento de Nonoai, localizado no extremo noroeste do Rio Grande do Sul, foi fundado em 1848 e posteriormente ratificado pela Lei de Terras. Sua criação inseriu-se em uma política do governo provincial que visava concentrar o maior número possível de indígenas em áreas reduzidas, buscando minimizar contatos violentos e garantir a segurança para a ocupação de brancos. Esse processo foi fundamental para a estratégia de expansão colonial do Império, servindo para liberar os territórios originais dos nativos para o assentamento de colonos alemães e italianos, que começaram a estabelecer municípios na região a partir de 1824 e 1875, respectivamente. Além de atuar como um núcleo de controle populacional, o aldeamento facilitava o desenvolvimento de infraestruturas, como a abertura de estradas ligando as Províncias do Rio Grande e de São Paulo, consolidando a expropriação das terras que pertenciam aos Kaingang desde tempos imemoriais.  

No século XIX, a convivência no aldeamento de Nonoai foi permeada por conflitos e invasões sistemáticas perpetradas por membros da sociedade envolvente que ignoravam os direitos territoriais indígenas. Um caso notório de abuso de poder foi o de Cipriano Rocha de Loures, o primeiro diretor do núcleo, que aproveitou sua posição para cercar três léguas das terras mais férteis do aldeamento em benefício próprio. Loures utilizava táticas de manipulação psicológica e ameaças de perseguição armada para forçar a saída dos nativos, chegando a destruir suas habitações para garantir a posse dos campos. Paralelamente, o avanço de posseiros como Clementino Pacheco, que se apoderou de terras reclamadas pelos indígenas sem que o governo tomasse as devidas providências, gerou hostilidades que culminaram em tragédia. A resistência Kaingang manifestou-se de forma violenta em 1856, quando o índio Pedro Nicofim executou Pacheco e outras cinco pessoas, evidenciando que o grupo reagia com vigor contra a expropriação de seus territórios sagrados e de seus meios de subsistência.

A atuação das lideranças Kaingang no século XIX, personificada por figuras como os Caciques Nonoai, Doble e Condá, revelou um protagonismo complexo que utilizou múltiplas táticas para garantir a sobrevivência e a manutenção de partes de seu território tradicional. O Cacique Nonoai destacou-se por sua liderança histórica na defesa das terras de seus antepassados, legitimando a posse dos campos por meio de vínculos de nascimento e ancestralidade. Simultaneamente, líderes como Vitorino Condá e Doble adotaram estratégias de colaboracionismo, atuando como mediadores e batedores junto às autoridades provinciais e colonizadores em troca de segurança, alimentos e gratificações para seus grupos. O Cacique Doble, de forma emblemática, exerceu um papel ambíguo: colaborava com o governo na captura de grupos arredios ao mesmo tempo em que participava de "correrias" (ataques violentos contra fazendas e lotes coloniais) para resistir à expropriação de suas áreas de caça e coleta. Dessa forma, a agência indígena manifestou-se na capacidade de alternar entre a diplomacia, muitas vezes mediada por missionários, e o uso da força física, transformando o aldeamento em um espaço de negociação e resistência estratégica frente ao avanço da sociedade expansionista.

O estudo conclui que os Kaingang não foram vítimas passivas do processo colonizador, mas sim sujeitos ativos que desenvolveram lógicas e comportamentos próprios diante das adversidades impostas pela expansão da sociedade nacional. Dentro da estrutura dos aldeamentos, esses grupos reelaboraram estratégias de sobrevivência, utilizando tais espaços não apenas como locais de confinamento, mas como uma forma tática de manter parte de seu território original e estabelecer um campo necessário de negociação com a sociedade envolvente para garantir sua existência física e cultural. A agência indígena manifestou-se por meio de ações diversificadas e complexas, que incluíram desde alianças estratégicas e colaboracionismo com autoridades provinciais até atos de resistência direta e violenta, demonstrando que os Kaingang lutaram ativamente pela preservação de seus interesses e terras tradicionais. Essa postura histórica de protagonismo fundamenta a relevância atual do grupo, que hoje constitui o povo indígena mais numeroso do Brasil Meridional e continua a mobilizar sua concepção tradicional de território para lutar pela demarcação e recuperação de suas terras ancestrais.

A sobrevivência cultural dos Kaingang evidencia que, apesar da expropriação brutal sofrida a partir de meados do século XIX, este povo permanece como o grupo indígena mais numeroso do Sul do Brasil e a maior representação da família linguística Macro-Jê. Com uma população estimada em aproximadamente 29 mil indivíduos distribuídos por dezenas de áreas entre São Paulo e o Rio Grande do Sul, o grupo transformou sua resiliência demográfica em um motor para novas reivindicações. Atualmente, a luta pela terra é encarada como a questão mais urgente para o grupo, uma vez que a sobrevivência física e a manutenção da identidade étnica dependem diretamente desse suporte territorial. Nesse contexto, os Kaingang buscam intensificar alianças institucionais e sociais para garantir a demarcação de espaços tradicionais ainda disponíveis, utilizando argumentos que unem a memória dos antepassados aos seus próprios significados socio-cosmológicos. Essa resistência contínua demonstra que, tal como no passado, a sociedade Kaingang segue reelaborando estratégias para garantir seu direito ao território e à posteridade.

\begin{thebibliography}{99}

\bibitem{curta2012}
	CURTA, Jussara Costa. \textbf{Continuidade e mudança do território Kaingang do Rio Grande do Sul: um estudo de caso do aldeamento de Nonoai}. 2012.

\bibitem{mabilde1983} 
MABILDE, Pierre F. A. B. \textbf{Apontamentos sobre os indígenas selvagens da Nação Coroados das matas da Província do Rio Grande do Sul}: 1836-1866. São Paulo: IBRASA, Fundação Nacional Pró-Memória, 1983.

\end{thebibliography}

\end{document}
